\answer{ex:07:1}{$V=\kappa p/\bigl(\kappa p+(1-\kappa)(1-p)\bigr)$: $0.059$, $0.20$, $0.50$; la médiane vaut $0$, alors que le point~\eqref{eq:expectile} varie continûment avec $\kappa$.}
\answer{ex:07:2}{$p=1/3$, $x=0.75$, $\sigma_r=0.35$, $V(0.2)=0.083$.}
\answer{ex:07:3}{$V=\sqrt\kappa/(\sqrt\kappa+\sqrt{1-\kappa})$, de $0.25$ à $0.75$; dispersion $\propto\sqrt\eta$; pour les gouttes $V=0.25+0.5\kappa$, les distributions diffèrent de $0.05$ aux neurones extrêmes, il faut $\eta\lesssim0.01$.}
\answer{ex:07:4}{(a)~$J^*=2\sqrt{C\bar R}$. (b)~Surcoût $\bigl(k^{1/2}+k^{-1/2}\bigr)/2-1$: $6\,\%$ pour $k=2$ et $25\,\%$ pour $k=4$; une précision \og à un facteur deux près\fg{} suffit.}
\answer{ex:07:5}{Pic d'aire $R(e^{-1}-e^{-2})\approx0.23R$, soit $2.3\,\text{s}$ de rampe; si la dopamine signale $V$, il n'y a pas de pic, le niveau est simplement multiplié par $e$ d'un seul coup.}
\answer{ex:07:6}{Un événement déplace le poids de $\propto A\tau_f/(\tau_e+\tau_f)$, le niveau donne une erreur $s\tau_f$: $9\,\text{s}\le\tau_f\le17\,\text{s}$.}
\answer{ex:07:7}{Au départ, $\Delta P=0.80$, $M=0.90$. (a)~$\Delta P=0.74$ ($-8\,\%$), $M=0.43$ ($-52\,\%$); (b)~$\Delta P=0.40$ ($-50\,\%$), $M=0.80$ ($-11\,\%$): double dissociation.}
\answer{ex:07:8}{$N_{\text{eff}}\to2+\rho^2N\approx10^6$ essais: $10^4$ fois moins qu'avec un seul fil, mais une fuite de $1\,\%$ coûte encore un million d'essais.}
